\documentclass[12pt]{article}
\usepackage[utf8]{inputenc}
\usepackage[T1]{fontenc}
\usepackage{hyperref}
\title{Yagua Language Resources for K-12 Curriculum Development}
\author{Compiled by Claude Sonnet 5.5 (claude-sonnet-5-5) for Andr\'as Kornai}
\date{}
\begin{document}
\maketitle

\section{Introduction}
Yagua (ISO 639-3: yad; Glottocode yagu1244) is a Peba-Yagua language of the Loreto region of Peru, along the Putumayo, Amazonas and Ampiyacu rivers. About 5,700 speakers are reported. Search results indicate a shift toward Spanish in some communities.

\section{Primary Educational Resources}
Search results state that Yagua is used as a medium of instruction in some primary schools in the Yagua-speaking region. No specific textbooks or primers were found, and current programme details were not confirmed.

\section{Linguistic and Reference Materials}
\begin{itemize}
\item Paul Powlison, \textit{Diccionario Yagua-Castellano} (1995, 771 pp.), published in Lima by the Ministerio de Educaci\'on and SIL. The catalogue entry is at \url{https://glottolog.org/resource/reference/id/40727}
\item Doris Payne, \textit{Aspects of the Grammar of Yagua: A Typological Perspective} (UCLA dissertation, 1985). \url{https://glottolog.org/resource/reference/id/305221}
\item Language Description Heritage entry for Yagua. \url{https://ldh.clld.org/languages/yad}
\item Yagua language overview. \url{https://en.wikipedia.org/wiki/Yagua_language}
\end{itemize}
A grammar sketch by Doris and Thomas Payne (1990) was mentioned in the results, but no URL was seen for it.

\section{Religious and Community Resources}
Nothing specific was found.

\section{Digital and Audio Resources}
\begin{itemize}
\item UNESCO World Atlas of Languages entry for Yagua in Peru. \url{https://ar.wal.unesco.org/countries/peru/languages/yagua}
\item Omniglot page on Yagua. \url{https://www.omniglot.com/writing/yagua.htm}
\end{itemize}

\section{Curriculum Development Context}
Yagua has a substantial dictionary and a full grammar, both published or deposited as scholarly works. It is reportedly used in primary schools. A curriculum could build on the Powlison dictionary and Payne grammar, ideally with Yagua teachers and communities. No ready-made K-12 materials were found.
\end{document}